\section{Introduction}\label{sec:introduction}

Every positive integer has a unique representation as a sum of
non-consecutive Fibonacci numbers \cite{Lekkerkerker1952,Zeckendorf1972}.
An arbitrary binary word $\omega\in\{0,1\}^m$ determines a Fibonacci-weighted
integer $N_m(\omega)=\sum_{j=1}^m\omega_jF_{j+1}$ that generally violates the
adjacency constraint; reducing modulo $F_{m+2}$ and inverting the standard
bijection $V_m\colon X_m\to\{0,\dots,F_{m+2}-1\}$ yields the finite-window
fold
\[
  \Fold_m(\omega):=V_m^{-1}\bigl(N_m(\omega)\bmod F_{m+2}\bigr).
\]
The fiber multiplicity $d_m(x):=|\Fold_m^{-1}(x)|$ records how many raw words
stabilize to a given visible word~$x$.

A natural question arises: what structure does this fiber spectrum carry, and
how does it connect to the classical theory of Fibonacci partitions?  We show
that the answer is highly structured.  After an explicit affine permutation
of residues, each fiber multiplicity equals a Fibonacci-lag discrete
derivative of the classical partition function.  This single identity converts
the fiber spectrum into a discrete pressure formalism whose exponential
constants are the Fibonacci partition constants of Sanna~\cite{Sanna2025}.  We
also record an audited finite arithmetic window for the corresponding
recurrence polynomials.

The paper makes three contributions.  First, it identifies the finite-window
fiber multiplicities with an explicit affine reindexing of Fibonacci partition
differences, so single fibers inherit the exact formulas of Chow and
Slattery~\cite{ChowSlattery2021}.  Second, it transfers the full moment
sequence to Sanna's partition constants and realizes those constants as Perron
roots of finite-state kernels, yielding algebraicity, convex pressure,
canonical concentration, and zero-temperature asymptotics.  Third, it records
an audited arithmetic window \(q=9,\dots,17\), in which the recurrence
polynomials are computationally certified to be irreducible with full
symmetric Galois groups, together with audited linear-disjointness and product
Chebotarev data for the block \(q\in\{12,13,14,15\}\).

\subsection{Main results}

Let $R(n)$ be the number of partitions of~$n$ into distinct Fibonacci numbers
(with $F_1=F_2=1$ not distinguished), set $R^\dagger(n):=R(n)+R(n-1)$, and
write $T_q^\dagger(N):=\sum_{n=0}^{N-1}R^\dagger(n)^q$.  For integers
$q\ge 1$, set $S_q(m):=\sum_{x\in X_m}d_m(x)^q$.  Also write
\[
  d_m^\#(r):=d_m(V_m^{-1}(r)),
\]
and let \(\pi_m\) denote the explicit affine permutation of
\(\ZZ/F_{m+2}\ZZ\) constructed in \Cref{sec:residue-affine}.

\begin{theoremA}
For every integer $n$ with $0\le n<F_{m+2}$,
\[
  d_m^\#(\pi_m(n))
  =
  R^\dagger(n)-R^\dagger(n-F_{m+1}).
\]
\end{theoremA}

\noindent
The fiber spectrum is the Fibonacci-lag discrete derivative of the indexed
partition function.  Exact Zeckendorf formulas of Chow and
Slattery~\cite{ChowSlattery2021} therefore apply directly to individual fibers.

\begin{theoremB}
For every integer $q\ge 1$,
\[
  T_q^\dagger(F_{m+1})\le S_q(m)\le T_q^\dagger(F_{m+2}),
\]
and consequently $S_q(m)\asymp_q \lambda_q^m$, where $\lambda_q$ is the
partition constant of Sanna~\cite[Thm.~1.1]{Sanna2025}.
\end{theoremB}

\begin{theoremC}
For $q\ge 2$, the Perron root of the finite-state collision kernel of
Section~\ref{sec:moment-kernel} equals $\lambda_q$.  More precisely,
$\lambda_q$ is the spectral radius of a nonnegative integer matrix of
dimension polynomial in~$q$; in particular, $\lambda_q$ is an algebraic
integer.  The pressure sequence $p_0:=\log\varphi$,
$p_q:=\log\lambda_q$ is convex, and its discrete slopes
$\Delta_q:=p_q-p_{q-1}$ satisfy $\Delta_q\nearrow\frac12\log\varphi$.
\end{theoremC}

\begin{theoremD}
For each fixed $q\ge 1$ and every $\varepsilon>0$, the $q$-canonical law
$\mu_{m,q}\propto d_m(x)^q$ is exponentially concentrated on the adjacent
pressure-slope band $[\Delta_q,\Delta_{q+1}]$:
\[
  \mu_{m,q}\bigl\{x:m^{-1}\log d_m(x)\notin
  [\Delta_q-\varepsilon,\,\Delta_{q+1}+\varepsilon]\bigr\}
  \le \exp\bigl(-(\varepsilon+o(1))m\bigr).
\]
\end{theoremD}

\begin{theoremE}
If
\[
  B_{m,q}(\varepsilon):=
  \left\{
    x\in X_m:
    \Delta_q-\varepsilon\le m^{-1}\log d_m(x)\le \Delta_{q+1}+\varepsilon
  \right\},
\]
then both its cardinality $|B_{m,q}(\varepsilon)|$ and its visible mass under
the law \(p_m(x):=d_m(x)/2^m\) satisfy two-sided exponential estimates with
exponents determined by $p_q,\Delta_q,\Delta_{q+1}$.
\end{theoremE}

\begin{theoremF}
Diverging tilts $q(m)\to\infty$ select the zero-temperature thickness
$\frac12\log\varphi$: if \(D_m:=\max_{x\in X_m}d_m(x)\), then
$D_m^{1/m}\to\sqrt{\varphi}$, and
$\log S_{q(m)}(m)=(\frac12\log\varphi+o(1))\,m\,q(m)$.
\end{theoremF}

\medskip\noindent
The escalation summarized by Theorems~A--F proceeds as follows.
Theorem~A identifies single fibers with partition differences.
Theorem~B transfers the moment growth to Fibonacci windows of Sanna's
power sums.  Theorem~C realizes the partition constants as finite-state
Perron roots and extracts a convex pressure function.  Theorems~D and~E
use the pressure geometry for canonical concentration and microcanonical
band bounds.  Theorem~F passes to the zero-temperature limit.

Beyond the thermodynamic formalism, the same partition constants carry
arithmetic content.  The quadratic case $q=2$ yields the cubic
$\lambda^3-2\lambda^2-2\lambda+2$, in agreement with Chow and
Jones~\cite{ChowJones2024}.  In the arithmetic window $q=9,\dots,17$, the
audited recurrence polynomials of $S_q(m)$ are computationally certified to be
irreducible over~$\QQ$ with full symmetric Galois groups on their roots
(Section~\ref{sec:chebotarev}); for $q\in\{12,13,14,15\}$ the same audit
indicates linear disjointness and the corresponding product Chebotarev
density.

\subsection{Related work}

Zeckendorf~\cite{Zeckendorf1972} and Lekkerkerker~\cite{Lekkerkerker1952}
established the uniqueness of Fibonacci representations.  Earlier analyses of
Fibonacci representations and their counting functions include
Carlitz~\cite{Carlitz1968}, Robbins~\cite{Robbins1996},
Weinstein~\cite{Weinstein2016}, and Berstel~\cite{Berstel2001}.  Chow and
Slattery~\cite{ChowSlattery2021} gave exact formulas for $R(n)$ via the
Zeckendorf expansion; our Theorem~A shows that finite-window fiber
multiplicities are, after an explicit affine permutation, the Fibonacci-lag
discrete derivatives of these same partition counts.  Chow and
Jones~\cite{ChowJones2024} identified the cubic governing the variance of
$R(n)$; our cubic recurrence for $S_2(m)$ gives a finite-window explanation.
Sanna~\cite[Thm.~1.1]{Sanna2025} proved that $\sum_{n<N}R(n)^q\asymp_q
N^{(\log\lambda_q)/\log\varphi}$; Theorem~B transfers these constants to the
collision moments.  The finite-state realization of
Section~\ref{sec:moment-kernel} draws on classical Perron--Frobenius theory
\cite{Seneta2006,LindMarcus1995} and on the automata-theoretic normalization
viewpoint developed by Frougny~\cite{Frougny1991FibonacciAutomata,Frougny1992NumbersAutomata}
and by Ahlbach, Usatine, Frougny, and
Pippenger~\cite{AhlbachUsatineFrougnyPippenger2013}.  For related entropy and
\(L^q\)-spectrum questions in the golden-ratio setting, see
Sidorov--Vershik~\cite{SidorovVershik1998} and
Lau--Ngai~\cite{LauNgai1998}.

\subsection{Organization}

Section~\ref{sec:preliminaries} defines the finite-window fold and records
its basic properties.  Section~\ref{sec:residue-affine} proves the affine
transfer principle and the partition-difference formula.
Section~\ref{sec:second-collision} establishes the Fibonacci-window moment
transfer and treats the explicit quadratic case.
Section~\ref{sec:moment-kernel} develops the finite-state realization, convex
pressure geometry, canonical selection, and zero-temperature concentration.
Section~\ref{sec:ledger} deduces the integer-order R\'enyi entropy rates of
the visible law.  Section~\ref{sec:chebotarev} records an audited finite
arithmetic window for the recurrence polynomials in the range
$q=9,\dots,17$.
Section~\ref{sec:conclusion} summarizes.  Appendix~\ref{app:transducer}
gives the corrected padded transducer construction for the collision kernels,
and Appendix~\ref{app:certification} records the exact certificate tables used
for the audited arithmetic window.
